\chapter*{ABSTRAK}
\addcontentsline{toc}{chapter}{ABSTRAK}

\begin{center}

    \textbf{Desain Arsitektur Transformasi Digital pada Kementerian Pertahanan} \\
    \vspace{0.4cm}
    oleh \\
    \vspace{0.4cm}
    Marcel Suandi Tambing\\
    23524316\\
\end{center}


\begin{spacing}{1.0}
Transformasi digital telah menjadi fondasi utama pertahanan modern yang
mengutamakan informasi dan pengambilan keputusan berbasis data. Namun,
implementasinya di Kementerian Pertahanan Republik Indonesia (Kemhan RI) masih
terkendala fragmentasi sistem, data silo, lemahnya interoperabilitas lintas unit dan
matra, serta kesenjangan kompetensi SDM digital, di tengah belum tersedianya
kerangka arsitektur yang terpadu. Penelitian ini bertujuan merancang arsitektur
transformasi digital Kemhan RI sebagai landasan menuju pertahanan cerdas
(\textit{smart defence}). Penelitian menggunakan \textit{Design Science Research Methodology}
(DSRM) yang dipadukan dengan TOGAF \textit{Architecture Development Method} (ADM).
Hasil analisis kondisi saat ini menunjukkan bahwa TOGAF ADM standar belum
mengakomodasi karakteristik pertahanan, khususnya aspek komando, kendali, dan
kesadaran situasional. Sebagai solusi, dirancang kerangka TOGAF Pertahanan
Indonesia yang menambahkan Fase C5ISR \textit{Architecture} sehingga terdiri atas sepuluh
fase, dilengkapi sembilan prinsip arsitektur dan usulan pembentukan unit
\textit{Digital Transformation Officer} (DTO). Artefak dievaluasi melalui
\textit{Enterprise Architecture Scorecard}\texttrademark{} dan oleh para \textit{Expert Judgement}.
Penelitian ini diharapkan menjadi pedoman strategis bagi Kemhan RI dalam
mewujudkan ekosistem digital pertahanan yang terintegrasi, aman, adaptif, dan
berkelanjutan.

\vspace{0.4cm}
\noindent\textbf{Kata kunci:} \textit{arsitektur transformasi digital, TOGAF ADM, DSRM, pertahanan cerdas, C5ISR.}
\end{spacing}